% math3.tex
%
% A superscript binds to exactly one token (or one braced group):
% x^2n is not "x to the 2n" -- it is "x to the 2" followed by a stray
% letter n, and it compiles just fine while meaning the wrong thing.
% Braces fix that: x^{2n} groups the whole exponent.
%
% Stacking a second caret on top, like x^2^3, is not "x to the (2 to
% the 3)" either -- TeX has no rule for which exponent binds first, so
% it refuses outright: "Double superscript." Group with braces instead.

% I AM NOT DONE

\documentclass{article}
\begin{document}

The expression $x^2n$ compiles, but it typesets the exponent as $2$
with a trailing $n$ -- not the exponent $2n$ you probably meant.
Chasing that mistake with more carets, an even worse attempt is
$x^2^3$.

\end{document}
